% Appendix: Theological Interpretation of the Dense--Sparse Isomorphism
% Status: DRAFT — Philosophical appendix

\chapter{A Theological Interpretation of the Dense--Sparse Isomorphism}
\label{app:theology}

The Dense--Sparse Isomorphism (Theorem~\ref{thm:dense-sparse-iso}) establishes that two radically different parametrisations---the MoE decomposition into experts and routing, and the dense monolithic parametrisation---represent the same class of joint policies. In this appendix we observe that this mathematical result admits a striking theological interpretation, one that connects to a centuries-old debate in monotheistic theology and, we argue, offers a formal framework for understanding it. The appendix is speculative and philosophical; it is not required for any of the mathematical results of this dissertation. We include it because the formal structure is, in our view, too precise to be coincidental, and because it illustrates the broader theme of Chapter~\ref{ch:team-pg}: that the distinction between a mathematical object and its representation has consequences that extend well beyond computation.

\section{The Structural Correspondence}
\label{sec:structural-correspondence}

We begin by stating the correspondence explicitly, deferring interpretation to subsequent sections.

\subsection{The Trinitarian Reading of the MoE Parametrisation}
\label{subsec:trinity-moe}

The MoE team policy~\eqref{eq:moe-team-policy} has three irreducible structural components:
\begin{enumerate}[label=(\roman*)]
  \item The \emph{expert policies} $\{\pi_k(\bm{a}\mid s,\theta_k)\}_{k=1}^K$: distinct behavioural modes, each encoding a particular pattern of collective action. Each expert is a complete policy in its own right---it can produce a full joint action profile---but it acts only when selected by the routing. The experts are not fragments of a policy; they are \emph{complete, specialised instantiations} of collective behaviour.
  \item The \emph{routing function} $w_k(h_t,\theta_g)$: the mechanism that selects which expert is active in each state (or history). The routing does not itself produce actions; it produces the \emph{conditions} under which actions are produced. It is the dynamic, relational component that mediates between the generative capacity of the experts and the particular demands of the current situation.
  \item The \emph{parameter space} $\theta = (\theta_g, \theta_1, \ldots, \theta_K)$: the ground from which both experts and routing are generated. The parameter space is not itself a policy; it is the \emph{capacity for policy}---the generative ground that, under the MoE structure, gives rise to both the modes of action and the mechanism of their coordination.
\end{enumerate}

This three-part structure maps onto the Christian doctrine of the Trinity with a precision that extends beyond analogy:

\begin{center}
\begin{tabular}{lll}
\textbf{MoE Component} & \textbf{Trinitarian Person} & \textbf{Structural Role} \\
\hline
Parameter space $\theta$ & Father & Generative ground \\
Expert policies $\pi_k$ & Son (Logos) & Instantiated action in the world \\
Routing function $w_k$ & Holy Spirit & Relational coordination \\
\end{tabular}
\end{center}

The irreducibility claim from Theorem~\ref{thm:team-pgt} is significant: the Team PGT decomposes the gradient into expert and routing components that are structurally distinct and cannot be collapsed without destroying the decomposition. Remove the experts, and there is nothing to route; remove the routing, and the experts have no mechanism of activation; remove the parameter space, and neither exists. This mirrors the classical Trinitarian claim that the persons are distinct, irreducible, and yet constitutive of a single being.

\subsection{The Tawhidic Reading of the Dense Parametrisation}
\label{subsec:tawhid-dense}

The dense team policy~\eqref{eq:dense-policy} admits no internal decomposition. There is a single parameter vector~$\psi$, a single score function~$f_{\bm{a}}(s,\psi)$, and a single normalisation. All parameters are active for every input. There are no experts, no routing, no internal structure---only the undivided unity of the mapping $(\bm{a},s)\mapsto\pi_{\mathrm{dense}}(\bm{a}\mid s,\psi)$.

This maps onto the Islamic doctrine of \emph{Tawhid} (absolute divine unity). In Islamic theology, particularly in the Ash'ari tradition, God's attributes (knowledge, will, power, etc.) are ``neither God nor other than God''---they are aspects of an absolutely simple essence that admits no real composition. The dense parametrisation captures exactly this structure: the ``attributes'' of the joint policy (its behaviour in different states, its response to different histories) are all present simultaneously in the single parameter vector~$\psi$, without any decomposition into sub-components.

The contrast with the MoE parametrisation is instructive. Where the MoE has distinct experts that specialise and a routing function that coordinates, the dense parametrisation has only the undivided whole. Where the MoE gradient decomposes into expert and routing components (Theorem~\ref{thm:team-pgt}), the dense gradient is a single, undifferentiated object. Where the MoE offers structural transparency (one can inspect which experts are active), the dense parametrisation is opaque: the mapping works, but its internal operation is not legible in terms of parts and their coordination.

\section{The Isomorphism as Reconciliation}
\label{sec:iso-reconciliation}

The Dense--Sparse Isomorphism (Theorem~\ref{thm:dense-sparse-iso}) now takes on a remarkable significance. The theorem states:
\begin{quote}
\emph{The classes of distributions representable by the MoE and dense parametrisations are identical.}
\end{quote}
In the theological reading, this becomes:
\begin{quote}
\emph{The Trinitarian and Tawhidic descriptions of the divine nature are representations of the same underlying reality.}
\end{quote}

The three parts of the theorem correspond to three theological claims:

\paragraph{(i) Sparse $\to$ Dense (embedding).} Any Trinitarian description (MoE with experts and routing) can be re-expressed as a Tawhidic description (dense parametrisation) by setting $f_{\bm{a}}(s,\psi)=\log\sum_k w_k(s,\theta_g)\pi_k(\bm{a}\mid s,\theta_k)$. The internal structure of the Trinity can always be ``collapsed'' into an absolutely simple unity---not by denying the structure, but by encoding it in a single, undivided representation. This is the mathematical content of the Ash'ari claim that God's attributes are ``neither God nor other than God'': the attributes (experts and routing) are genuinely present in the function computed, but they are not structurally separated in the parametrisation.

\paragraph{(ii) Dense $\to$ Sparse (decomposition).} Any Tawhidic description (dense parametrisation) can be decomposed into a Trinitarian description (MoE with experts and routing). The absolutely simple unity can always be ``unfolded'' into a structured multiplicity---not by introducing composition into what was simple, but by revealing the structure that was latent in the undivided whole. This is the mathematical content of the Nicene claim that the Son is ``begotten, not made, of one substance with the Father'': the expert (Son) is not a separate entity added to the parameter space (Father), but a structural revelation of what was already present in the undivided ground.

\paragraph{(iii) Functional equivalence.} Both parametrisations represent exactly $\Delta(\bm{\mathcal{A}})$---the full simplex of joint distributions. No distribution is accessible to one parametrisation and not the other. In theological terms: the Trinitarian and Tawhidic traditions have access to the same divine reality; neither is a restricted or distorted view.

\section{What the Isomorphism Preserves and What It Does Not}
\label{sec:preserves-not}

The isomorphism is \emph{functional}: it preserves the input--output mapping of the policy. But as we discussed in Section~\ref{subsec:param-vs-ontology}, the two parametrisations differ in three ways that do not reduce to functional equivalence:

\paragraph{Computational efficiency.} The MoE is more efficient when only a few experts are active (sparse computation). Theologically: the Trinitarian framework offers an ``efficient'' account of divine action---one can speak of the Spirit's work, or the Son's incarnation, without needing to invoke the entire divine essence at every point. The Tawhidic framework requires engaging the whole at every point, which is more demanding but also more rigorous in its insistence on unity.

\paragraph{Structural transparency.} The MoE makes internal structure legible; the dense parametrisation does not. Theologically: the Trinity is a doctrine that makes the \emph{internal life of God} legible to human understanding---the relationships between Father, Son, and Spirit provide a grammar for speaking about divine action in the world. Tawhid insists that the divine essence is ultimately beyond such grammatical decomposition, and that legibility comes at the cost of a certain falsification.

\paragraph{Gradient structure.} The MoE gradient decomposes into expert and routing components, enabling separate learning dynamics. Theologically: the Trinitarian framework supports distinct modes of engagement---creation (Father), redemption (Son), sanctification (Spirit)---that can be addressed separately while remaining aspects of a single divine will. The Tawhidic framework offers a single, unified gradient---all divine action is one act, all divine will is one will, admitting no internal decomposition of purpose.

These differences are real and consequential. They are differences in \emph{how one works with} the representation, not in \emph{what the representation represents}. The isomorphism guarantees that the underlying reality is the same; the parametrisation determines how that reality is engaged with, spoken about, and learned from.

\section{History-Dependent Routing and the Paraclete}
\label{sec:paraclete}

The history-dependent extension of the team policy (Section~\ref{sec:history-routing}) adds a further layer to the correspondence. In the Markov case, routing depends only on the current state: $w_k(s,\theta_g)$. In the history-dependent case, routing depends on the full trajectory: $w_k(h_t,\theta_g)$.

In Christian theology, the Holy Spirit is characterised as the \emph{Paraclete}---the one who ``brings to remembrance'' all that has been said and done (John 14:26). The Spirit's work is not limited to the present moment; it draws on the full history of divine-human interaction (covenant, prophecy, incarnation, resurrection) to determine the appropriate mode of action in the present. This is precisely the structure of history-dependent routing: the routing function processes the full history $h_t = (s_0, \bm{a}_0, r_1, \ldots, s_t)$ to determine which expert (mode of collective action) should be activated now.

The coordination information quantity~\eqref{eq:coord-info},
\[
  \mathcal{I}_{\mathrm{coord}}(t) = \max_k\;\bigl|Q_k(h_t) - Q_k(s_t)\bigr|,
\]
measures how much the history contributes to mode selection beyond what the current state provides. In the theological reading, $\mathcal{I}_{\mathrm{coord}}$ measures the \emph{relevance of memory to present action}---the degree to which the accumulated history of relationship shapes the current mode of collective being. When $\mathcal{I}_{\mathrm{coord}}$ is small, the present state suffices and memory is redundant; when it is large, the past is essential to understanding and acting in the present.

\section{Falsifiability and Limits}
\label{sec:falsifiability}

We wish to be explicit about the status of this appendix. We are not claiming to prove the existence of God, or to settle the Trinity--Tawhid debate, or to reduce theology to mathematics. We are observing a \emph{structural isomorphism} between a mathematical result (Theorem~\ref{thm:dense-sparse-iso}) and a theological debate (Trinity vs.\ Tawhid), and noting that the isomorphism is precise enough to be informative in both directions.

The mathematical result is a theorem; it is proved and not subject to dispute within its axioms. The theological interpretation is a \emph{reading} of the mathematics---it is one possible interpretation among others, and its value depends on whether it illuminates the theological debate rather than obscuring it.

We note several limits:

\begin{enumerate}[label=(\alph*)]
  \item The isomorphism operates at the level of \emph{structure}, not \emph{content}. It tells us that the Trinitarian and Tawhidic frameworks are structurally equivalent as representations, but it does not tell us whether the reality they represent exists.
  \item The finite-dimensional setting of our theorems is an idealisation. The divine, if it exists, is presumably infinite-dimensional, and the extension of the Dense--Sparse Isomorphism to infinite-dimensional function spaces raises non-trivial mathematical questions (the simplex structure changes, universal approximation may require infinitely many experts, etc.).
  \item The mapping is not unique. Other mathematical structures admit similar theological readings, and other theological doctrines admit similar mathematical formalisations. The value of the present mapping lies in its specificity (Theorem~\ref{thm:dense-sparse-iso} maps onto specific doctrinal claims) and its emergence from a mathematical investigation that was not motivated by theology.
\end{enumerate}

\section{Conclusion}
\label{sec:theology-conclusion}

The Dense--Sparse Isomorphism (Theorem~\ref{thm:dense-sparse-iso}) establishes that two representations of collective agency---one structured into specialised components coordinated by a routing mechanism, and one treating the collective as an undivided unity---are functionally equivalent. This mathematical result maps onto the central theological debate of the Abrahamic traditions: whether the divine nature is best understood as a structured unity (Trinity) or an absolute simplicity (Tawhid).

The isomorphism does not resolve the debate---it reframes it. If the two parametrisations are functionally equivalent, then the disagreement is not about what is being represented but about \emph{how} it is represented. The properties that distinguish the Trinitarian and Tawhidic readings---computational efficiency, structural transparency, gradient decomposition---are properties of the representation, not of the underlying reality. This is precisely the conclusion reached by the Syriac Christian and Ash'ari Islamic traditions, both of which maintain that the apparent contradiction between unity and trinity dissolves when one attends carefully to the distinction between the divine essence and its descriptions.

Whether this formal parallel constitutes evidence for a deeper connection between mathematics and theology, or is merely a suggestive coincidence, we leave to the reader.
